\section{Алгебраически замкнутые поля}

Пусть $F$~--- поле. Говорят, что многочлен \emph{раскладывается} в $F[X]$ (на линейные множители), если он является
произведением многочленов степени не выше $1$ из $F[X]$.

\begin{proposition}[1.43]
Для поля $\Om$ следующие условия равносильны:
\begin{enumerate}[label=(\alph*),leftmargin=*]
\item каждый непостоянный многочлен из $\Om[X]$ раскладывается в $\Om[X]$ на линейные множители;
\item каждый непостоянный многочлен из $\Om[X]$ имеет в $\Om$ хотя бы один корень;
\item неприводимые многочлены в $\Om[X]$~--- это в точности многочлены степени $1$;
\item каждое конечное расширение $\Om$ совпадает с $\Om$.
\end{enumerate}
\end{proposition}

\begin{proof}
(a)$\Rightarrow$(b). Если непостоянный $f$ раскладывается, то среди его множителей есть линейный $c(X-a)$, и $a$~---
корень $f$.

(b)$\Rightarrow$(c). Многочлены степени $1$ неприводимы всегда. Если $\deg f>1$, то по (b) у $f$ есть корень $a\in\Om$, и
$f=(X-a)g$ с $\deg g=\deg f-1\ge1$ (\rf{1.7}), т.\,е.\ $f$ приводим.

(c)$\Rightarrow$(a). Индукция по $n=\deg f\ge1$. При $n=1$ многочлен $f$ уже линеен. Пусть $n>1$. По (c) многочлен $f$
приводим: $f=gh$ с $1\le\deg g,\deg h<n$. По предположению индукции $g$ и $h$ раскладываются на линейные множители, а
значит, и $f$. (Иначе: $\Om[X]$~--- факториальное кольцо, и $f$~--- произведение неприводимых, т.\,е.\ по (c) линейных,
многочленов.)

(c)$\Rightarrow$(d). Пусть $E$~--- конечное расширение $\Om$ и $\alpha\in E$. По \rf{1.30} элемент $\alpha$ алгебраичен над
$\Om$; его минимальный многочлен неприводим, поэтому по (c) имеет степень $1$, т.\,е.\ равен $X-\alpha$, и $\alpha\in\Om$.
Итак, $E=\Om$.

(d)$\Rightarrow$(c). Пусть $f\in\Om[X]$ неприводим. Тогда $\Om[X]/(f)$~--- расширение $\Om$ степени $\deg f$ (\rf{1.25}), и
по (d) оно совпадает с $\Om$, т.\,е.\ $\deg f=1$.
\end{proof}

В условии (d) слово <<конечное>> можно заменить на <<алгебраическое>>: если $E/\Om$ алгебраично и $\alpha\in E$, то
$\Om[\alpha]$~--- конечное расширение $\Om$ (\rf{1.27c}), и по (d) $\alpha\in\Om$. Таким образом, алгебраически замкнутое
поле не имеет собственных алгебраических расширений.

\begin{definition}[1.44]
\begin{enumerate}[label=(\alph*),leftmargin=*]
\item Поле $\Om$ называется \emph{алгебраически замкнутым}, если оно удовлетворяет равносильным условиям \rf{1.43}.
\item Поле $\Om$ называется \emph{алгебраическим замыканием} своего подполя $F$, если $\Om$ алгебраически замкнуто и
алгебраично над $F$.
\end{enumerate}
\end{definition}

Например, основная теорема алгебры (\rf{5.6}) утверждает, что поле $\CC$ алгебраически замкнуто. Оно является
алгебраическим замыканием $\R$ (каждое $a+bi\in\CC$~--- корень многочлена $X^2-2aX+a^2+b^2\in\R[X]$). Но $\CC$ не является
алгебраическим замыканием $\Q$: оно не алгебраично над $\Q$ (\rf{1.33}, \rf{1.34}). Конечное поле не бывает алгебраически
замкнутым: многочлен $1+\prod_{a\in F}(X-a)$ не имеет корней в конечном поле $F$.

\begin{proposition}[1.45]
Пусть поле $\Om$ алгебраично над подполем $F$ и каждый многочлен $f\in F[X]$ раскладывается в $\Om[X]$ на линейные
множители. Тогда $\Om$ алгебраически замкнуто (и, следовательно, является алгебраическим замыканием $F$).
\end{proposition}

\begin{proof}
Проверим условие \rf{1.43}(b). Пусть $f=a_nX^n+\dots+a_0\in\Om[X]$, $a_i\in\Om$, $n\ge1$. Возьмём неприводимый делитель $g$
многочлена $f$ в $\Om[X]$; в поле $\Om'=\Om[X]/(g)$, конечном над $\Om$ (\rf{1.25}), у $g$, а значит и у $f$, есть корень
$\alpha$. Рассмотрим поля
\[
F\ \subset\ F[a_0,\dots,a_n]\ \subset\ F[a_0,\dots,a_n,\alpha]\ \subset\ \Om'.
\]
Элементы $a_0,\dots,a_n\in\Om$ алгебраичны над $F$ по условию, поэтому $F[a_0,\dots,a_n]$~--- поле, конечное над $F$
(\rf{1.30} и замечание после него). Элемент $\alpha$ алгебраичен над $F[a_0,\dots,a_n]$~--- он корень многочлена $f$ с
коэффициентами из этого поля, поэтому по тому же замечанию и $F[a_0,\dots,a_n,\alpha]$~--- поле, конечное над
$F[a_0,\dots,a_n]$. По \rf{1.20} поле $F[a_0,\dots,a_n,\alpha]$ конечно над $F$, и по \rf{1.30} $\alpha$ алгебраичен
над $F$, т.\,е.\ является корнем ненулевого многочлена $h\in F[X]$. По условию $h=c\prod_j(X-\beta_j)$ с $\beta_j\in\Om$.
Подставляя $\alpha$ в это равенство в поле $\Om'$, получаем $\prod_j(\alpha-\beta_j)=0$, т.\,е.\ $\alpha=\beta_j\in\Om$ для
некоторого $j$. Итак, $f$ имеет корень в $\Om$.
\end{proof}

\par\noindent\begin{minipage}{\linewidth}
\[
\begin{tikzpicture}[baseline]
\node (F) at (0,0) {$F$};
\node (A) at (0,1.3) {$F[a_0,\dots,a_n]$};
\node (B) at (0,2.6) {$F[a_0,\dots,a_n,\alpha]$};
\node (O) at (3.6,1.3) {$\Om$};
\node (O1) at (3.6,2.6) {$\Om'=\Om[X]/(g)$};
\draw (F)--(A)--(B);
\draw (A)--(O)--(O1);
\draw (B)--(O1);
\draw[decorate,decoration={brace,amplitude=5pt},cnew] (-1.4,0.15)--(-1.4,2.45) node[midway,left=5pt]{конечно};
\end{tikzpicture}
\]
\figcap{Корень $\alpha$ многочлена $f\in\Om[X]$ найден в конечном расширении $\Om'$; так как башня слева конечна, $\alpha$
алгебраичен над $F$, а все корни многочленов из $F[X]$ уже лежат в $\Om$.}
\end{minipage}\par\medskip

На самом деле достаточно требовать, чтобы каждый непостоянный $f\in F[X]$ имел в $\Om$ хотя бы \emph{один} корень (\rf{6.5}).

\begin{proposition}[1.46]
Пусть $F$~--- поле, а $\Om$~--- область целостности, содержащая $F$ в качестве подкольца. Тогда множество
\[
\overline F=\{\alpha\in\Om\mid \alpha\ \text{алгебраичен над}\ F\}
\]
является полем. Оно называется \emph{алгебраическим замыканием $F$ в $\Om$}.
\end{proposition}

\begin{proof}
Пусть $\alpha,\beta\in\overline F$. Подкольцо $F[\alpha,\beta]=F[\alpha][\beta]$ области $\Om$~--- область целостности,
порождённая над $F$ двумя алгебраическими элементами; по \rf{1.32} это поле, алгебраическое над $F$ (и даже конечное: \rf{1.30} и замечание после него).
Значит, $F[\alpha,\beta]\subset\overline F$. В частности, $\alpha\pm\beta$, $\alpha\beta$ и (при $\beta\ne0$) $\alpha/\beta$
лежат в $\overline F$; ясно также, что $F\subset\overline F$.
\end{proof}

\begin{corollary}[1.47]
Пусть $\Om$~--- алгебраически замкнутое поле. Для любого его подполя $F$ алгебраическое замыкание $E$ поля $F$ в $\Om$
является алгебраическим замыканием $F$.
\end{corollary}

\begin{proof}
Поле $E$ (\rf{1.46}) алгебраично над $F$ по определению. Каждый многочлен $f\in F[X]$ раскладывается в $\Om[X]$ на линейные
множители $X-\alpha_i$; все его корни $\alpha_i$ алгебраичны над $F$ и потому лежат в $E$, т.\,е.\ $f$ раскладывается уже в
$E[X]$. По \rf{1.45} $E$ алгебраически замкнуто.
\end{proof}

Алгебраическая замкнутость $\Om$ здесь существенна. Например, алгебраическое замыкание $\Q$ в $\Om=\R$~--- поле
$\R\cap\Qal$ вещественных алгебраических чисел (здесь $\Qal$~--- поле всех алгебраических чисел, см.\ \rf{1.47a})~---
алгебраично над $\Q$, но не алгебраически замкнуто: у $X^2+1$ нет в нём
корней.

Итак, если принять основную теорему алгебры (\rf{5.6}), у каждого подполя $F\subset\CC$ есть алгебраическое замыкание~---
алгебраическое замыкание $F$ в $\CC$. Это замыкание даже \emph{каноническое}: оно задано без произвольного выбора как
конкретное подмножество $\CC$; что означает это слово и почему для абстрактного поля такого выбора нет, обсуждается в
\rf{6.8a}. В главе 6 с помощью леммы Цорна (или принципа ультрафильтра) доказывается, что алгебраическое замыкание есть у
любого поля (раздел <<Существование алгебраического замыкания>> главы~6) и что оно единственно с точностью до
(неединственного) изоморфизма (\rf{6.6}).

\begin{example}[1.47a: поле алгебраических чисел]
Обозначим через $\Qal$ алгебраическое замыкание $\Q$ в $\CC$, т.\,е.\ поле всех алгебраических чисел. По \rf{1.47} это
алгебраическое замыкание $\Q$, и по \rf{1.33} оно счётно.

\emph{Расширение $\Qal/\Q$ алгебраично, но бесконечно.} Для каждого $n\ge1$ многочлен $X^n-2$ неприводим над $\Q$
(\rf{1.16a}(a)), поэтому $[\Q(\sqrt[n]2):\Q]=n$ (\rf{1.27c}). Если бы $[\Qal:\Q]=N<\infty$, то из
башни $\Q\subset\Q(\sqrt[n]2)\subset\Qal$ по \rf{1.20} следовало бы $n\mid N$ для всех $n$, что неверно, например, при
$n=N+1$. Итак, $[\Qal:\Q]=\infty$.

Это показывает, что <<алгебраическое>> и <<конечное>>~--- разные понятия: конечное расширение алгебраично (\rf{1.30}), но
алгебраическое расширение конечно лишь тогда, когда оно порождено конечным числом элементов. Каждый элемент
$\alpha\in\Qal$ лежит в конечном подрасширении $\Q(\alpha)$, и $\Qal$~--- объединение этих подрасширений, но их степени
не ограничены в совокупности.
\end{example}

\begin{aside}[1.47b]
Бесконечность $[\Qal:\Q]$ следует и из теоремы Артина--Шрайера: если алгебраически замкнутое поле $C$ конечно над
подполем $K\ne C$, то $C=K(\sqrt{-1})$ и $[C:K]=2$ (как для $\CC\supset\R$). Если бы $\Qal$ было конечно над $\Q$, то
$\Qal=\Q(i)$. Но $\sqrt2\notin\Q(i)$: иначе подполе $\Q(\sqrt2)$ степени $2$ совпало бы с $\Q(i)$ (по \rf{1.20} промежуточное
поле степени $2$ в расширении степени $2$ совпадает со всем полем), а $\Q(\sqrt2)\subset\R$, тогда как $i\notin\R$.
Конечно, это стрельба из пушки по воробьям: элементарное рассуждение \rf{1.47a} гораздо проще.
\end{aside}

\begin{aside}[1.47c: историческое замечание]
Отдельные классы полей (числовые поля, поля функций) изучались и раньше, но первое систематическое изложение теории
абстрактных полей дал Э.~Штейниц (1910). Он ввёл понятие простого подполя,
различил сепарабельные и несепарабельные расширения, показал, что каждое поле получается алгебраическим расширением
чисто трансцендентного расширения, и доказал, что у каждого поля есть алгебраическое замыкание, единственное с точностью
до изоморфизма. Его работа сильно повлияла на следующее поколение алгебраистов (Э.~Нётер, Б.~Л.~ван дер Варден, Э.~Артин).
\end{aside}
